\documentclass[10pt,a4paper]{amsart}
\usepackage[margin=19mm]{geometry}
\usepackage{amsmath,amssymb,mathtools}
\usepackage[scr=rsfs,cal=euler]{mathalfa}
\usepackage{xfrac}
\usepackage[unicode,hidelinks,bookmarks=false]{hyperref}
\newcommand{\ie}{i.e.\ }
\newcommand{\cf}{cf.\ }
\newcommand{\resp}{respectively\ }
\newcommand{\Subsection}{\subsection}
\providecommand{\nobreakdash}{\nobreak\hbox{-}}
\setlength{\emergencystretch}{2em}
\multlinegap0pt
\usepackage{bm}
\usepackage{bbm}
\usepackage{dsfont}
\usepackage{xspace}
\usepackage{cancel}
\usepackage{mathtools}
\usepackage[shortlabels]{enumitem}
\usepackage{amsthm}
\makeatletter
\@ifundefined{theorem}{\newtheorem{theorem}{Theorem}}{}
\@ifundefined{lemma}{\newtheorem{lemma}[theorem]{Lemma}}{}
\@ifundefined{proposition}{\newtheorem{proposition}[theorem]{Proposition}}{}
\makeatother
\newcommand{\II}{\mathcal{I}}
\newcommand{\N}{\ensuremath{\mathbb{N}}}
\newcommand{\al}{\alpha}
\newcommand{\f}{\infty}
\newcommand{\lle}{\preccurlyeq}
\newcommand{\set}[1]{\left\{#1\right\}}
\newcommand{\si}{\sigma}
\newcommand{\w}{\mathbf{w}}
\renewcommand{\a}{ {\mathbf{a}}}
\renewcommand{\u}{\mathbf{u}}
\renewcommand{\v}{\mathbf{v}}
\title[Entropy plateaus, transitivity and bifurcation sets for the ]{Entropy plateaus, transitivity and bifurcation sets for the $\beta$-transformation with a hole at $0$}
\author{Pieter Allaart, Derong Kong}
\date{Source: arXiv:2304.06892v3 (CC BY 4.0)}
\begin{document}
\maketitle

\noindent\emph{Source and licence.} Entropy plateaus, transitivity and bifurcation sets for the $\beta$-transformation with a hole at $0$. Version arXiv:2304.06892v3 (CC BY 4.0), distributed under the Creative Commons Attribution 4.0 International licence, as recorded in the arXiv licence metadata for this exact version and shown on the abstract page https://arxiv.org/abs/2304.06892v3. Full rights notice: licenses/arxiv-2304.06892-CC-BY-4.0.txt.\par\smallskip

\noindent\textit{Editorial scope.} Counted unit: the Proposition whose source label is \texttt{prop:non-transitivity-intervals} in the article (source0.tex lines 2812--2821, source label \texttt{prop:non-transitivity-intervals}), delivered together with the article's own complete original proof of it (source0.tex lines 2823--2908, one \texttt{proof} environment of 86 lines). No mathematics is written, repaired or re-derived: statement and proof are the article's own text line for line. Required context retained uncounted: lem:quasi-greedy expansion-alpha-q (lines 446--452); lem:repeating-preserves-order (lines 2347--2349); lem:always-below-alpha (lines 2555--2560); eq:always-below-alpha (lines 2557--2559); lem:v\_k-decreasing (lines 2594--2600); lem:header-suffix (lines 2617--2619); lem:v\_k-upper-bound (lines 2626--2637); lem:v\_k-star (lines 2715--2728). Every definition, display and label of those lines is retained unchanged. Omitted from this package and not redistributed: 1--445, 453--2346, 2350--2554, 2560--2593, 2601--2616, 2620--2625, 2638--2714, 2729--2811, 2822--2822, 2909--4006. Nothing inside a retained excerpt is omitted or altered: every retained body is delivered exactly as the article's lines are written, so no omission receipt of any kind accompanies this package. Citations: the retained excerpts carry no citation command at all, so no bibliography entry of the article is transcribed and no reference is redistributed. Reference adaptation: none was needed and none was made; every reference inside the delivered unit resolves inside it, so no unresolved reference or citation is produced. Printed slips kept literal: no printed word, symbol, hypothesis or number of the counted statement or of its proof was altered. Layout and class: the article's own class is \texttt{amsart} with packages including hyperref, a4wide, amsmath, inputenc, amssymb, amsopn, epsfig, amsfonts, latexsym, graphicx, enumerate, xcolor, mathrsfs, tikz; the wrapper is the lane's pinned \texttt{amsart} editorial wrapper at 10pt on A4, which restates the theorem-like environments the retained text needs and adds the article's own macro definitions that the retained text needs (\texttt{\textbackslash II}, \texttt{\textbackslash N}, \texttt{\textbackslash a}, \texttt{\textbackslash al}, \texttt{\textbackslash f}, \texttt{\textbackslash lle}, \texttt{\textbackslash set}, \texttt{\textbackslash si}, \texttt{\textbackslash u}, \texttt{\textbackslash v}, \texttt{\textbackslash w}), with extra wrapper packages (bm, bbm, dsfont, xspace, cancel, mathtools, enumitem). The delivered numbering is the unit's own. Assets: no retained span contains a figure environment, a graphics-inclusion command, a PDF-page inclusion, a table or a TikZ picture; no image file is copied into this package, the package declares no asset dependency and no image file is redistributed with it. Licence: version arXiv:2304.06892v3 of this article is distributed under the Creative Commons Attribution 4.0 International licence, as recorded in the arXiv licence metadata for this exact version and shown on the abstract page https://arxiv.org/abs/2304.06892v3. A full rights notice is delivered at licenses/arxiv-2304.06892-CC-BY-4.0.txt. Characters: every retained excerpt is pure ASCII, so the delivered engine, XeLaTeX with its default Latin Modern text font, reports no missing character.
\begin{lemma} \label{lem:quasi-greedy expansion-alpha-q}
{Let $k\in\N$. The restriction of the map $\beta\mapsto\al(\beta)$ to $(k,k+1]$ is an increasing bijection from $\beta\in (k,k+1]$ to the set of sequences $\a=\al_1\al_2\dots\in\set{0,1,\dots,k}^\N$ such that $\al_1=k$ and}
%The map $\beta\mapsto \al(\beta)$ is strictly increasing and bijective from $(1,2]$ to the set of sequences $(\al_i)\in\{0,1\}^\N$ satisfying
\[
0^\f\prec\si^n((\al_i))\lle (\al_i)\quad\forall n\ge 1.
\]
\end{lemma}

\begin{lemma} \label{lem:repeating-preserves-order}
If $\v$ and $\w$ are Lyndon words and $\v\prec\w$, then $\v^\f\prec \w^\f$.
\end{lemma}

\begin{lemma} \label{lem:always-below-alpha}
For any $k$ and any two integers $0\leq i_1<i_2\leq l_k$,
\begin{equation} \label{eq:always-below-alpha}
v_{k,i_1+1}\dots v_{k,i_2}\lle \al_1\dots\al_{i_2-i_1}.
\end{equation}
\end{lemma}

\begin{lemma} \label{lem:v_k-decreasing}
Let $k\geq 2$, and suppose $\v_{k}\neq \v_{k-1}$. Then
\begin{enumerate}[{\rm(i)}]
\item $\v_{k-1}$ is not a prefix of $\v_{k}$; and
\item $\v_k \prec \v_{k-1}$.
\end{enumerate}
\end{lemma}

\begin{lemma} \label{lem:header-suffix}
For each $k\in\N$, the word $\al_1\dots\al_{j_k}$ can occur in $\v_k$ at most once, and only at the end. In other words, the truncated word $v_{k,1}\dots v_{k,l_k-1}$ does not contain the word $\al_1\dots\al_{j_k}$.
\end{lemma}

\begin{lemma} \label{lem:v_k-upper-bound}
\begin{enumerate}[{\rm(i)}]
\item For each $k\in\N$, we have
\begin{equation} \label{eq:v_k-upper-bound}
\v_k\lle\al_{m+1}\dots \al_{m+l_k} \qquad\forall\, 0\leq m<j_k.
\end{equation}
\item Suppose either $k=1$, or $k\geq 2$ and $\v_k\neq \v_{k-1}$. If $\v_k$ is not $\beta$-Lyndon, then
\begin{equation} \label{eq:strict-v_k-upper-bound}
\v_k\prec\al_{m+1}\dots \al_{m+l_k} \qquad\forall\, 0\leq m<j_k.
\end{equation}
\end{enumerate}
\end{lemma}

\begin{lemma} \label{lem:v_k-star}
Let $k\geq 1$, and suppose $\v_k$ is not $\beta$-Lyndon. If $k\geq 2$ and $\v_k=\v_{k-1}$, then $\v_k^*=\v_{k-1}^*$. Otherwise, there exist a nonnegative integer $r$ and a word $\u$ not ending in $\al_1\dots\al_{j_k}^-$ such that
\begin{equation} \label{eq:v_k-expression}
\v_k=\u(\al_1\dots\al_{j_k}^-)^r\al_1\dots\al_{j_k},
\end{equation}
and we have
\begin{equation} \label{eq:v_k-star}
\v_k^*=\begin{cases}
\v_{k-1}^* & \mbox{if $k\geq 2$ and $\u=\v_{k-1}^-$},\\
\u^+ & \mbox{otherwise}.
\end{cases}
\end{equation}
%On the other hand, if $\v_k$ does not end in $\al_1\dots\al_{j_k}$, then $\v_k^*=\v_k$.
\end{lemma}

\begin{proposition} \label{prop:non-transitivity-intervals}
Let $k,\ell\in\N$ with $k<\ell$, and suppose $I_k=[\underline{t}_k,\overline{t}_k)$ and $I_\ell=[\underline{t}_\ell,\overline{t}_\ell)$ belong to $\II$. Then:
\begin{enumerate}[{\rm(i)}]
\item Either $\overline{t}_{\ell}\leq\underline{t}_k$ or $I_{\ell}\subseteq I_k$. In other words, $I_\ell$ is either contained in $I_k$ or else lies completely to the left of $I_k$.
\item $I_{\ell}\subseteq I_k$ if and only if
\begin{equation} \label{eq:containment-condition}
\v_{\ell}=\v_k\quad\mbox{or}\quad \v_{\ell}\ \mbox{begins with}\ \v_k^-(\al_1\dots\al_{j_k}^-)^n\al_1\dots\al_{j_k}\ \mbox{for some $n\geq 0$}.
\end{equation}
\end{enumerate}
\end{proposition}

\begin{proof}
It suffices to show that \eqref{eq:containment-condition} implies $I_{\ell}\subseteq I_k$, and $\overline{t}_{\ell}>\underline{t}_k$ implies \eqref{eq:containment-condition}.

Assume first that \eqref{eq:containment-condition} holds. If $\v_{\ell}=\v_k$, then $\underline{t}_{\ell}>\underline{t}_k$ and $\overline{t}_{\ell}=\overline{t}_k$, so $I_{\ell}\subseteq I_k$. If $\v_{\ell}\neq \v_k$, then by iterating Lemma \ref{lem:v_k-decreasing}, $\v_{\ell}\prec \v_k$ and $\v_{\ell}$ does not extend $\v_k$. Hence, we also have $\v_{\ell}^*\lle \v_k^*$ since the map $\v\mapsto \v^*$ which sends a Lyndon word $\v$ to the smallest $\beta$-Lyndon word $\v^*$ greater than or equal to $\v$, is clearly nondecreasing. Therefore, {by Lemma \ref{lem:repeating-preserves-order}}, $\overline{t}_{\ell}\leq \overline{t}_k$. On the other hand, $\v_\ell$ begins with $\v_k^-(\al_1\dots\al_{j_k}^-)^n\al_1\dots\al_{j_k}$ for some $n\geq 0$, and this implies $\underline{t}_\ell\geq \underline{t}_k$. Thus, $I_\ell\subseteq I_k$.

Next, we show by induction that for all $\ell\geq k$,
$\overline{t}_{\ell}>\underline{t}_k$ implies \eqref{eq:containment-condition}.
This is trivial for $\ell=k$, so take $\ell'>k$ and assume that $\overline{t}_{\ell}>\underline{t}_k$ implies \eqref{eq:containment-condition} for all $k\leq \ell<\ell'$. Suppose $\overline{t}_{\ell'}>\underline{t}_k$, so
\begin{equation} \label{eq:v_l-v_k-inequality}
(\v_{\ell'}^*)^\f \succ \v_k^-(\al_1\dots\al_{j_k}^-)^\f.
\end{equation}
Assume $\v_{\ell'}\neq \v_k$, as otherwise there is nothing to show.
We consider two cases:

\medskip
{\em Case 1.} $\v_{\ell'}$ is $\beta$-Lyndon. Then $\v_{\ell'}^*=\v_{\ell'}$ and hence
\[
\v_{\ell'}^\f \succ \v_k^-(\al_1\dots\al_{j_k}^-)^\f.
\]
Since $\v_k$ and $\v_{\ell'}$ are both Lyndon, $\v_{\ell'}^\f$ cannot begin with $\v_k$, and hence it must begin with $\v_k^-$. Thus, in view of the constraint \eqref{eq:always-below-alpha}, there exists some integer $r\geq 0$ such that $\v_{\ell'}^\f$ begins with
\begin{equation} \label{eq:v_l-beginning}
\v_k^-(\al_1\dots\al_{j_k}^-)^r\al_1\dots\al_{j_k}.
\end{equation}
Suppose, by way of contradiction, that $\v_{\ell'}$ is a proper prefix of this last word. {If $\v_{\ell'}$ is a prefix of $\v_k^-(\al_1\dots\al_{j_k}^-)^r$, then $\si^{|\v_{\ell'}|}(\v_{\ell'}^\f)=\v_{\ell'}^\f$ begins both with \eqref{eq:v_l-beginning}, and with a proper suffix of \eqref{eq:v_l-beginning} which contains the word $\al_1\dots\al_{j_k}$. Hence, there is an earlier occurrence of $\al_1\dots\al_{j_k}$ in \eqref{eq:v_l-beginning}. But this is impossible by Lemmas \ref{lem:quasi-greedy expansion-alpha-q}, \ref{lem:always-below-alpha} and \ref{lem:header-suffix}. This leaves the case when
\[
\v_{\ell'}=\v_k^-(\al_1\dots\al_{j_k}^-)^r \al_1\dots\al_m \qquad\mbox{for some $1\leq m<j_k$}.
\]
We then have (by shifting \eqref{eq:v_l-beginning} to the left by $|\v_{\ell'}|$)
\[
\big(\v_k^-(\al_1\dots\al_{j_k}^-)\big)_{1:j_k-m}=\al_{m+1}\dots\al_{j_k}.
\]
This implies that $l_k=|\v_k|>j_k-m$ by Lemma \ref{lem:v_k-upper-bound}, and 
\[
v_{k,1}\dots v_{k,j_k-m}=\al_{m+1}\dots\al_{j_k}.
\]
Again using Lemma \ref{lem:v_k-upper-bound}, this forces
\[
v_{k,j_k-m+1}\dots v_{k,l_k}\lle \al_{j_k+1}\dots\al_{l_k+m}=v_{k,1}\dots v_{k,l_k+m-j_k},
\]
}
contradicting that $\v_k$ is Lyndon. Therefore, $\v_{\ell'}$ begins with $\v_k^-(\al_1\dots\al_{j_k}^-)^r\al_1\dots\al_{j_k}$.

% so
%\[
%\v_{\ell'}=\v_k^-(\al_1\dots\al_{j_k}^-)^p \al_1\dots\al_m \qquad\mbox{for some $0\leq p\leq r$ and $1\leq m<j_k$}.
%\]
%{[We still need to explain why $\v_{\ell'}$ can't be a proper prefix of $\v_k^-$.]}
%(Note $\v_{\ell'}$ cannot end in $\al_1\dots\al_{j_k}^-$ because $\v_{\ell'}$ is Lyndon.) Then it is easy to see that $\v_{\ell'}$ does not contain the word $\al_1\dots\al_{j_k}$ (this follows from Lemma \ref{lem:v_k-upper-bound}). However, $\v_{\ell'}^\f$ does contain $\al_1\dots\al_{j_k}$ because it begins with \eqref{eq:v_l-beginning}. Therefore, {in view of Lemma \ref{lem:quasi-greedy expansion-alpha-q}} it must be that $\al_1\dots\al_m \v_k^-\al_1\dots\al_{j_k}^-$ contains $\al_1\dots\al_{j_k}$, and this word must begin inside the block $\al_1\dots\al_m$. So let $i<m$ be such that $\al_{i+1}\dots\al_m \v_k^-\al_1\dots\al_{j_k}^-$ begins with $\al_1\dots\al_{j_k}$. Then necessarily $j_k-(m-i)<|\v_k|=l_k$ in view of Lemma \ref{lem:v_k-upper-bound}. But then we have
%\[
%v_{k,1}\dots v_{k,j_k-(m-i)}=\al_{m-i+1}\dots\al_{j_k},
%\]
%so again using Lemma \ref{lem:v_k-upper-bound} it follows that
%\[
%v_{k,j_k-(m-i)+1}\dots v_{k,l_k}\lle \al_{j_k+1}\dots\al_{m-i+l_k}=v_{k,1}\dots v_{k,l_k+m-i-j_k},
%\]
%contradicting that $\v_k$ is Lyndon. Therefore, $\v_{\ell'}$ begins with $\v_k^-(\al_1\dots\al_{j_k}^-)^r\al_1\dots\al_{j_k}$.

\medskip
{\em Case 2.} $\v_{\ell'}$ is not $\beta$-Lyndon. Let $\ell\leq \ell'$ be the largest integer such that $\v_{\ell}=\v_{\ell'}$.
Then $\v_{\ell}\neq \v_{\ell-1}$, $\v_{\ell}\neq \v_k$ since $\v_{\ell'}\neq \v_k$, and of course $\v_{\ell}$ is not $\beta$-Lyndon. By Lemma \ref{lem:v_k-star},
\begin{equation} \label{eq:v-ell-form}
\v_{\ell}=\u(\al_1\dots\al_{j_{\ell}}^-)^r\al_1\dots\al_{j_{\ell}}
\end{equation}
for some word $\u$ not ending in $\al_1\dots\al_{j_{\ell}}^-$ and some $r\geq 0$. If $\u=\v_k^-$, then $\v_\ell$ extends $\v_k^-\al_1\dots\al_{j_k}$ {since $j_\ell>j_k$}. So assume $\u\neq\v_k^-$. Then there are two possibilities:
%by Lemma \ref{lem:v_k-star} there are two possibilities:

(a) $\u=\v_{\ell-1}^-$. Then $\v_{\ell'}^*=\v_\ell^*=\v_{\ell-1}^*$, so $(\v_{\ell-1}^*)^\f \succ \v_k^-(\al_1\dots\al_{j_k}^-)^\f$ by \eqref{eq:v_l-v_k-inequality}. This means $\overline{t}_{\ell-1}>\underline{t}_k$, so by the induction hypothesis, $\v_{\ell-1}$ begins with $\v_k^-(\al_1\dots\al_{j_k}^-)^n\al_1\dots\al_{j_k}$ for some $n\geq 0$, since $\v_{\ell-1}=\u^+\neq \v_k$. Now either (i) $\v_{\ell-1}=\v_k^-(\al_1\dots\al_{j_k}^-)^n\al_1\dots\al_{j_k}$, in which case {\eqref{eq:v-ell-form} implies}
\[
\v_\ell=\v_{\ell-1}^-(\al_1\dots\al_{j_{\ell}}^-)^r\al_1\dots\al_{j_{\ell}}=\v_k^-(\al_1\dots\al_{j_k}^-)^{n+1}(\al_1\dots\al_{j_{\ell}}^-)^r\al_1\dots\al_{j_{\ell}},
\]
and so $\v_\ell$ begins with $\v_k^-(\al_1\dots\al_{j_k}^-)^{n+1}\al_1\dots\al_{j_k}$ because $j_\ell>j_k$; or (ii) $\v_{\ell-1}$ properly extends $\v_k^-(\al_1\dots\al_{j_k}^-)^n\al_1\dots\al_{j_k}$, in which case $\u$, and hence $\v_\ell$, begins with ${\v_k^-}(\al_1\dots\al_{j_k}^-)^n\al_1\dots\al_{j_k}$.

(b) $\u\neq\v_{\ell-1}^-$. Then $\v_{\ell'}^*=\v_\ell^*=\u^+$, and hence,
\[
(\u^+)^\f\succ \v_k^-(\al_1\dots\al_{j_k}^-)^\f.
\]
This implies that $(\u^+)^\f$ begins with $\v_k^-(\al_1\dots\al_{j_k}^-)^{q}\al_1\dots\al_{j_k}$ for some ${q}\geq 0$. The same argument as in Case 1 then shows that $\u^+$ {begins with} $\v_k^-(\al_1\dots\al_{j_k}^-)^{q}\al_1\dots\al_{j_k}$. If $\u^+$ properly extends this word, then so does $\u$ and hence so does $\v_\ell$.
Otherwise, {\eqref{eq:v-ell-form} yields
\[
\v_\ell=\u(\al_1\dots\al_{j_{\ell}}^-)^r\al_1\dots\al_{j_\ell}=\v_k^-(\al_1\dots\al_{j_k}^-)^{q+1}(\al_1\dots\al_{j_{\ell}}^-)^r\al_1\dots\al_{j_\ell},
\]
}
so that $\v_\ell$ {begins with} $\v_k^-(\al_1\dots\al_{j_k}^-)^{{q}+1}\al_1\dots\al_{j_k}$, because $j_\ell>j_k$.

In both case (a) and case (b) we have shown that $\v_\ell$ {begins with} $\v_k^-(\al_1\dots\al_{j_k}^-)^n\al_1\dots\al_{j_k}$ for some $n\geq 0$. {Since $\v_{\ell'}=\v_\ell$}, this completes the proof.
\end{proof}

\end{document}
